\section{Related Work}
\label{sec:related-work}

The review supports two design decisions. Section~\ref{sec:rw-disfluency}
examines what NLP systems remove, preserve or recover when processing
self-repair. Section~\ref{sec:rw-representation-choice} then asks which
semantic accounts can encode the missing distinction between withdrawn and
accepted content. Readers familiar with disfluency detection may begin with
the second subsection; readers familiar with semantic representation can use
the first to locate the project's NLP task.

\subsection{Disfluency Detection and Downstream Processing}
\label{sec:rw-disfluency}

Mainstream NLP approaches use \emph{disfluency} as an umbrella term for a
range of phenomena that interrupt the fluent flow of speech, including
self-repair
\citep{shribergPreliminariesTheorySpeech1994,senSemanticParsingDisfluent2021}.
The studies reviewed below differ in their coverage of these phenomena, but
typically address them within a shared processing task. Their outputs range
from cleaned text and annotated transcripts to meaning representations for
downstream applications
\citep{johnsonTAGbasedNoisychannelModel2004,futamiStreamingJointSpeech2023,senSemanticParsingDisfluent2021}.

In a task-oriented NLP pipeline, a downstream component that books a ticket
or retrieves information typically only needs the user's final intention: which
ticket to book or which information to return. When that is the output
requirement, an upstream disfluency detector can mark material for omission
from subsequent interpretation. Recovering the user's
final intention is therefore an appropriate objective for these immediate
downstream tasks. Applied to Figure~\ref{fig:siri-repair}, however, this
objective yields a request for Dublin's weather without recording that Berlin
was first proposed and then withdrawn. The following discussion traces this
representational issue through correction models, transcript annotation and
semantic parsing from disfluent input.

Early computational models used the structural correspondences between
reparandum and repair described by \citet{leveltMonitoringSelfrepairSpeech1983}
to guide detection and correction. \citet{heemanDetectingCorrectingSpeech1994},
for example, used them to detect self-repairs as speech unfolded and identify
material to remove. Noisy-channel approaches modelled the reparandum as
material introduced into a fluent source and recovered that source for
subsequent language understanding
\citep{johnsonTAGbasedNoisychannelModel2004,honalCorrectionDisfluenciesSpontaneous2003}.
Across these approaches, structural correspondence guides the recovery of
fluent text.

Later neural methods reduced the need for hand-designed matching features.
For example, \citet{zayatsRobustCrossdomainDisfluency2018} learned similarity
patterns useful for identifying reparanda, while retaining detection as the
output task. The scope of the output labels, however, had already received
attention. \citet{zayatsMultidomainDisfluencyRepair2014} criticised the focus
on reparandum detection without type distinctions. Their model distinguished
repetitions from other types of disfluency and added labels for repair spans,
providing a richer account of self-repair in the transcript.
\citet{futamiStreamingJointSpeech2023} likewise preserved disfluent material
in joint speech recognition and disfluency detection, using either span
boundary markers or token-level fluent/disfluent labels. These approaches
preserve information that deletion would remove, with varying detail about
disfluency types and spans. Their labels do not, however, explicitly
represent how the repair replaces the reparandum's contribution to the
utterance's meaning.

An alternative is to let language models handle disfluency implicitly,
without a separate detection stage.
\citet{senSemanticParsingDisfluent2021} used a sequence-to-sequence parser
with a pre-trained RoBERTa encoder to map disfluent queries directly to
the intent-and-slot representations of their fluent counterparts. When
trained only on fluent data, it achieved 2\% exact-match accuracy on 56
disfluent queries, compared with 81\% on their fluent counterparts.
Adding targeted disfluency training data improved performance. More recent
end-to-end voice systems integrate spoken input, tool calls and spoken
responses, likewise handling disfluency without a separate detection
output. Yet in Full-Duplex-Bench-v3 \citep{linFullDuplexBenchv3BenchmarkingTool2026}, the highest task-completion rate on
scenarios containing self-repair was 58.8\%, achieved by GPT-Realtime. Reported errors include
premature commitment to tool arguments, allowing reparanda to enter
API calls. These findings indicate that reliable disfluency handling
remains a challenge across both pre-trained semantic parsers and
end-to-end voice systems.

The approaches reviewed above offer ways to detect disfluency and recover
the speaker's final intention, but leave the semantic status of withdrawn
content largely outside their output requirements. Together with the
continuing contextual role of the reparandum discussed in the introduction,
this motivates exploring alternative ways of representing self-repair.

\subsection{Formal and Computational Semantic Accounts of RSR}
\label{sec:rw-representation-choice}

This subsection compares semantic accounts of RSR and related forms of
revision, focusing on whether their representations can serve as targets for
LLM-based semantic parsing.

Dynamic Syntax with Type Theory with Records (DS-TTR) builds meaning word by
word through DS parsing actions and expresses the resulting content as TTR
record types \citep{purverIncrementalSemanticConstruction2011}.
\citet{houghProcessingSelfrepairsIncremental2012} handle RSR by
returning to an earlier parse state and continuing along a new path.
Their processing record links words, parsing actions and semantic
representations, keeping the reparandum accessible. The final accepted TTR
record type alone, however, does not identify the withdrawn path and its
replacement; that information remains in the processing history.

This project requires the final meaning representation itself to retain the
reparandum, repair and replacement structure for inspection, comparison and
scoring. DS-TTR is therefore not adopted as the basis for the target
representation because it would shift the learning and evaluation target from
a final repair-aware semantic structure to a grammar-constrained derivation.

AMR represents sentence meaning as a graph
\citep{banarescuAbstractMeaningRepresentation2013} and can be linearised for
generation by pre-trained sequence-to-sequence models
\citep{fanLeveragingAMRGraph2024}, making it a potential candidate for this project. A
close precedent comes from work on disrupted sentences.
\citet{addleseeUnderstandingDisruptedSentences2023} use underspecified AMR
graphs to represent sentences with missing content and complete them using
follow-up input. Their task recovers missing content, whereas RSR replaces
content already expressed. In Figure~\ref{fig:siri-repair}, a repair-aware
representation must therefore retain both locations and mark the status of
each.

Standard AMR nevertheless requires extension on two points. First, its
limited treatment of scope \citep{bosSquibExpressivePower2016} lacks the
separate semantic contexts needed to \emph{quarantine} the reparandum.
Its content must be retained without entering the accepted meaning. Second,
AMR abstracts away from surface word order
\citep{banarescuAbstractMeaningRepresentation2013}, so a graph and its
linearisation do not record the order of replacement updates. This order matters
because reversing which expression replaces the other changes the accepted
meaning. An AMR extension would therefore need an explicit directed relation
or a constrained linearisation. These requirements favour a context-based
framework that already separates semantic contexts and relates successive
updates.

Within such a framework, correction must distinguish the content that is
given up from the content that replaces it.
\citet{vanleusenIncompatibilityContextDiagnosis2004} characterises this
relation through contextual incompatibility. Accepting the new proposal
requires giving up the content being revised, whether that content was stated
explicitly or inferred. In Figure~\ref{fig:siri-repair}, this means withdrawing
Berlin as the location and accepting Dublin while the surrounding request
remains unchanged.

Segmented Discourse Representation Theory (SDRT)
\citep{asherLogicsConversation2003} encodes connections between discourse
units through rhetorical relations. \citet{lascaridesAgreementDisputesCommitments2009}
use \textsc{Correction} to link the corrected unit to the correcting one,
representing the incompatibility condition proposed by
\citet{vanleusenIncompatibilityContextDiagnosis2004} as the denial of prior
content and the acceptance of its correction. Their account principally concerns
corrections across conversational turns. SBN inherits this theoretical
machinery by adopting SDRT's discourse-relation inventory to connect semantic
contexts \citep{bosSequenceNotationCatching2023}. The remaining task is to
specify how the reparandum and repair connect to the shared semantic structure
of one unfinished clause. Section~\ref{sec:formalism} develops this
configuration and its interpretation.
